\section{Problem 3: Debye Solid}

In Problem 2, you computed the thermodynamics of an Einstein solid by assuming that every oscillator has the same angular frequency \(\omega_E\). In this problem, you will replace the single Einstein frequency with a continuum of vibrational frequencies.

The Debye model describes a solid using a phonon density of states. Instead of assigning every vibrational mode the same frequency, the Debye model distributes the modes over frequencies from \(0\) to the Debye cutoff frequency \(\omega_D\).

The computational chain is
\[
\omega
\longrightarrow
g(\omega)
\longrightarrow
\langle E(\omega,T)\rangle
\longrightarrow
U(T)
\longrightarrow
C_V(T).
\]

The purpose of this problem is to show how a distribution of microscopic vibrational modes determines the macroscopic heat capacity of a solid.

\subsection{Physical model}

A crystalline solid with \(N_{\mathrm{atom}}\) atoms has
\[
N_{\mathrm{mode}}=3N_{\mathrm{atom}}
\]
vibrational modes.

In the Debye model, the allowed vibrational frequencies lie in the interval \(0\le \omega\le\omega_D\). The three-dimensional Debye density of states is
\[
g(\omega)=\frac{9N_{\mathrm{atom}}}{\omega_D^3}\omega^2,
\qquad
0\le\omega\le\omega_D.
\]

This density of states is normalized so that
\[
\int_0^{\omega_D}g(\omega)\,d\omega=3N_{\mathrm{atom}}.
\]
This condition means that the density of states contains the correct total number of vibrational modes.

Each mode is still a quantum harmonic oscillator. A mode with angular frequency \(\omega\) has average energy
\[
\langle E(\omega,T)\rangle
=
\frac12\hbar\omega
+
\frac{\hbar\omega}{
\exp\left(\frac{\hbar\omega}{k_{\mathrm B}T}\right)-1
}.
\]
The first term is the zero-point energy. The second term is the temperature-dependent thermal energy.

The total internal energy is obtained by integrating over all phonon frequencies:
\[
U(T)=
\int_0^{\omega_D}
g(\omega)\langle E(\omega,T)\rangle\,d\omega.
\]

The heat capacity is
\[
C_V(T)=\frac{dU}{dT}.
\]
In your implementation, compute \(C_V(T)\) using a centered finite difference:
\[
C_V(T)\approx
\frac{U(T+\Delta T)-U(T-\Delta T)}{2\Delta T}.
\]

\subsection{Numerical model}

Complete the provided skeleton implementation in \texttt{debye.py}. The constructor should store the physical parameters and create a frequency grid on the interval \(0\le\omega\le\omega_D\).

Use a uniform grid \(\omega_i=i\Delta\omega\), where \(i=0,1,\ldots,N_\omega-1\). The grid should cover the full interval from \(0\) to \(\omega_D\). Use this grid to approximate the internal energy integral by the finite sum
\[
U(T)\approx
\sum_i
g(\omega_i)
\langle E(\omega_i,T)\rangle
\Delta\omega.
\]

The method \texttt{frequency\_grid} should return the frequency grid and the spacing \(\Delta\omega\). The method \texttt{density\_of\_states} should evaluate \(g(\omega)\). The method \texttt{mode\_average\_energy} should evaluate \(\langle E(\omega,T)\rangle\). The method \texttt{internal\_energy} should approximate the Debye energy integral. The method \texttt{heat\_capacity} should compute \(C_V(T)\) from a finite difference of \(U(T)\).

Near \(\omega=0\), the factor
\[
\frac{\hbar\omega}{
\exp\left(\frac{\hbar\omega}{k_{\mathrm B}T}\right)-1
}
\]
has a finite limit. Your numerical implementation should handle the \(\omega=0\) point carefully so that it does not produce division by zero or \texttt{nan} values.

\subsection{Numerical checks}

First, check the normalization of the density of states. Numerically approximate
\[
\int_0^{\omega_D}g(\omega)\,d\omega
\]
using your frequency grid. The result should be close to \(3N_{\mathrm{atom}}\).

Second, check that the internal energy is positive:
\[
U(T)>0.
\]

Third, compute \(C_V(T)\) over a temperature grid. The heat capacity should be small at low temperature and should approach the high-temperature classical limit
\[
C_V\rightarrow3N_{\mathrm{atom}}k_{\mathrm B}.
\]

Fourth, check the effect of the number of frequency grid points \(N_\omega\). If \(N_\omega\) is too small, the integral over the density of states will be inaccurate. Repeat the calculation for several values of \(N_\omega\) and show that the density-of-states normalization and heat-capacity curve converge.

\subsection{Numerical experiments}

Use your implementation to study how different frequency ranges contribute to the heat capacity.

For each temperature \(T_j\), compute the contribution of each frequency grid point to the internal energy:
\[
u_i(T_j)=g(\omega_i)\langle E(\omega_i,T_j)\rangle\Delta\omega.
\]
Store these values in a two-dimensional array. Use this array to make a heat map showing which phonon frequencies contribute most strongly to the internal energy at different temperatures.

Then compute the cumulative mode count
\[
M(\omega_m)=\sum_{i=0}^{m}g(\omega_i)\Delta\omega.
\]
This should increase from \(0\) to approximately \(3N_{\mathrm{atom}}\). Plot \(M(\omega)\) versus \(\omega\) to show how the Debye density of states accumulates vibrational modes.

Also compute the cumulative internal-energy contribution
\[
U(\omega_m,T)=\sum_{i=0}^{m}g(\omega_i)\langle E(\omega_i,T)\rangle\Delta\omega.
\]
Use this to show how low-frequency and high-frequency modes contribute at different temperatures.

Finally, compare the numerical heat capacity against the expected limiting behavior. At low temperature, the Debye model predicts a \(T^3\)-type heat capacity. At high temperature, the heat capacity should approach \(3N_{\mathrm{atom}}k_{\mathrm B}\). Your results should show both trends if the frequency grid and temperature range are chosen well.

\subsection{Required plots}

Include the following plots:

\begin{enumerate}
    \item Debye density of states \(g(\omega)\) versus \(\omega\);
    \item cumulative mode count \(M(\omega)\) versus \(\omega\);
    \item internal energy \(U(T)\) versus temperature;
    \item heat capacity \(C_V(T)\) versus temperature;
    \item normalized heat capacity \(C_V/(3N_{\mathrm{atom}}k_{\mathrm B})\) versus temperature;
    \item heat map of frequency-resolved internal-energy contributions \(u_i(T_j)\);
    \item heat-capacity convergence for several values of \(N_\omega\).
\end{enumerate}

Each plot should have labeled axes and a short caption or explanation.

\subsection{Results}

Your results should report the numerical density-of-states normalization
\[
\sum_i g(\omega_i)\Delta\omega
\]
and compare it with \(3N_{\mathrm{atom}}\).

Report the internal energy and heat capacity over the chosen temperature range. State whether \(C_V(T)\) approaches the classical high-temperature limit. Also state whether the low-temperature behavior is consistent with the Debye model.

Report how the results change as \(N_\omega\) increases. The density-of-states normalization and heat-capacity curve should become more stable as the frequency grid is refined.

\subsection{Discussion}

In your discussion, explain how the Debye model improves on the Einstein model by using a distribution of vibrational frequencies instead of a single oscillator frequency.

Explain the meaning of the density of states \(g(\omega)\). The density of states tells how many vibrational modes occur in a small frequency interval near \(\omega\). Because \(g(\omega)\propto\omega^2\), the number of modes grows rapidly at higher frequencies.

Discuss the heat-capacity curve. At low temperature, only low-frequency modes are significantly thermally excited. At high temperature, all modes become thermally active and the heat capacity approaches \(3N_{\mathrm{atom}}k_{\mathrm B}\).

Discuss the frequency-resolved heat map. Explain which frequencies contribute at low temperature and how the contributing frequency range changes as temperature increases.

Discuss numerical convergence. Explain why the density-of-states normalization and heat capacity depend on the frequency grid resolution. State how you decided that your chosen \(N_\omega\) was large enough.

Conclude by summarizing the computational chain
\[
\omega
\longrightarrow
g(\omega)
\longrightarrow
\langle E(\omega,T)\rangle
\longrightarrow
U(T)
\longrightarrow
C_V(T).
\]
This chain shows how the Debye density of states connects microscopic vibrational modes to macroscopic heat capacity.

\subsection{Collaboration reminder}

No points will be deducted for appropriate collaboration, provided that the collaboration is disclosed and the submitted code, calculations, plots, explanations, and conclusions are your own.